% relatedwork.tex -- Related Work

Our work lies at the intersection of associative memory, energy-based modeling, and Langevin sampling.

The classical Hopfield network~\cite{hopfieldNeuralNetworksPhysical1982} stores binary patterns as local minima of a quadratic energy; retrieval proceeds by gradient descent but the storage capacity scales only as $\mathcal{O}(d)$~\cite{amitStoringInfiniteNumbers1985}. Dense associative memories~\cite{krotovDenseAssociativeMemory2016,krotovLargeAssociativeMemory2021} replace the quadratic interaction with higher-order polynomials, breaking the linear barrier. At the top of this hierarchy, the log-sum-exp energy of modern Hopfield networks~\cite{ramsauerHopfieldNetworksAll2021} achieves exponential capacity and is equivalent to transformer attention: each head performs one step of gradient descent on the energy. All of these works focus on \emph{retrieval}; sampling from the Boltzmann distribution of the same energy is not explored. The retrieval--generation duality dates to the Boltzmann machine~\cite{ackeleyLearningAlgorithmBoltzmann1985}, which shares the Hopfield energy but samples via Gibbs dynamics over binary states. Our formulation lifts this duality to the continuous modern Hopfield setting where Langevin dynamics scales naturally to high dimensions.

Energy-based models (EBMs) define $p(\mathbf{x})\propto\exp(-E(\mathbf{x}))$ with a neural-network energy trained by contrastive divergence~\cite{hintonTrainingProductsExperts2002}, score matching~\cite{hyvarinen2005estimation}, or noise contrastive estimation~\cite{gutmannHyvaerinenNoiseContrastive2010,songHowTrainEnergyBased2021}. Because the energy is a black box, its gradient must be learned. By contrast, the Hopfield energy has a closed-form gradient (identity minus softmax attention), requires no training, and provides analytic guarantees (smoothness, Lipschitz gradients, quadratic confinement). The Energy Transformer~\cite{hooverEnergyTransformer2024} brings Hopfield energies into a deep architecture for discriminative tasks but does not sample from the Boltzmann distribution. Score-based generative models~\cite{songGenerativeModelingEstimating2019,songScoreBasedGenerativeModeling2021} and diffusion models~\cite{hoDenoisingDiffusionProbabilistic2020,sohlDicksteinDeepUnsupervisedLearning2015} achieve remarkable sample quality but rely on a learned score network; our score is exact and analytic, at the cost of being restricted to a fixed memory $\mathbf{X}$. Normalizing flows~\cite{rezendeMohamed2015} construct generative models via learned invertible transformations but, like diffusion models, require training data and a learned bijection without exposing a closed-form score over a fixed memory. Deng et al.~\cite{dengLatentAlignment2018} introduce stochasticity into attention by replacing softmax with a latent variable trained via variational inference to model alignment uncertainty; our stochasticity arises from Langevin dynamics on a fixed energy with no variational objective.
